% GENERATED FILE. Edit canonical lessons, then run npm run generate:books.
% canonical-source-hash: fnv1a64:c5b3019879ca3f3e
% canonical-lessons: GU-C227-ija, GU-C227-gha, GU-C227-hadku, GU-C227-mackodavu, GU-C227-langdavu

\chapter{A fall in the street}
\label{ch:gu-a-fall-in-the-street}
\hlchaptermodality{227}

\begin{chapteropening}
\textbf{By the end of this chapter:} \emph{I can say that I am hurt and ask about a sprain.}

Complete the last lesson of chapter 227: I can say that I am hurt and ask about a sprain.
\end{chapteropening}

\section[\texorpdfstring{\gu{ઈજા} (ījā)}{ઈજા (ījā)}]{\gu{ઈજા} (ījā) — injury}

\label{lesson:GU-C227-ija}



\begin{quote}
Before the new one: say the Gujarati for to swallow, then the Gujarati for lung.
\end{quote}

\subsection*{You'll want to know: \gu{ઈજા}}
\textbf{\gu{ઈજા}} — \emph{ījā} — \textquotedblleft{}injury\textquotedblright{}.

\begin{quote}
\textbf{\gu{શું} \gu{તમને} \gu{ઈજા} \gu{થઈ}?} — \emph{śũ tamne ījā thaī?} — \textquotedblleft{}Did you get hurt?\textquotedblright{}
\end{quote}

\begin{cousinweb}
\textbf{\gu{ઈજા}} (\emph{ījā}) is a loan from Arabic \emph{īdhā'}, \textquotedblleft{}harm, hurt\textquotedblright{}, which reached Gujarati through Persian. Its Arabic root is \emph{'-dh-y} (Lane).
\end{cousinweb}

\begin{grammarlens}[title={What you've built}]
One more word for everyday life at A2.
\end{grammarlens}

\subsection*{Your turn}
\begin{itemize}
\raggedright
  \item \emph{Say it:} \emph{ījā}
  \item \emph{Say it:} \emph{ījā}, once more
  \item \emph{Say it:} \emph{śũ tamne ījā thaī?} — Did you get hurt?
  \item \emph{Recall:} say the Gujarati for to swallow, then the Gujarati for lung, then say \emph{ījā} again
\end{itemize}

\begin{tcolorbox}[breakable,colback=teal!4,colframe=teal!35!black,title={Before you move on}]
What does \gu{ઈજા} mean? (\textquotedblleft{}Injury\textquotedblright{}.) Say it once more.
\end{tcolorbox}
\section[\texorpdfstring{\gu{ઘા} (ghā)}{ઘા (ghā)}]{\gu{ઘા} (ghā) — wound, cut}

\label{lesson:GU-C227-gha}



\begin{quote}
Before the new one: say the Gujarati for lung, then the Gujarati for injury.
\end{quote}

\subsection*{You'll want to know: \gu{ઘા}}
\textbf{\gu{ઘા}} — \emph{ghā} — \textquotedblleft{}wound, cut\textquotedblright{}.

\begin{quote}
\textbf{\gu{આ} \gu{ઘા} \gu{મોટો} \gu{નથી}.} — \emph{ā ghā moṭo nathī} — \textquotedblleft{}This wound is not big.\textquotedblright{}
\end{quote}

\begin{cousinweb}
\textbf{\gu{ઘા}} (\emph{ghā}) is a tadbhava. It comes from Sanskrit \emph{ghāta}, \textquotedblleft{}killing; blow, bruise\textquotedblright{}, and CDIAL glosses the Gujarati word \textquotedblleft{}stroke, cut\textquotedblright{} (CDIAL 4460).
\end{cousinweb}

\begin{grammarlens}[title={What you've built}]
One more word for everyday life at A2.
\end{grammarlens}

\subsection*{Your turn}
\begin{itemize}
\raggedright
  \item \emph{Say it:} \emph{ghā}
  \item \emph{Say it:} \emph{ghā}, once more
  \item \emph{Say it:} \emph{ā ghā moṭo nathī} — This wound is not big.
  \item \emph{Recall:} say the Gujarati for lung, then the Gujarati for injury, then say \emph{ghā} again
\end{itemize}

\begin{tcolorbox}[breakable,colback=teal!4,colframe=teal!35!black,title={Before you move on}]
What does \gu{ઘા} mean? (\textquotedblleft{}Wound, cut\textquotedblright{}.) Say it once more.
\end{tcolorbox}
\section[\texorpdfstring{\gu{હાડકું} (hāḍkũ)}{હાડકું (hāḍkũ)}]{\gu{હાડકું} (hāḍkũ) — bone}

\label{lesson:GU-C227-hadku}



\begin{quote}
Before the new one: say the Gujarati for injury, then the Gujarati for wound.
\end{quote}

\subsection*{You'll want to know: \gu{હાડકું}}
\textbf{\gu{હાડકું}} — \emph{hāḍkũ} — \textquotedblleft{}bone\textquotedblright{}.

\begin{quote}
\textbf{\gu{આ} \gu{હાડકું} \gu{નબળું} \gu{નથી}.} — \emph{ā hāḍkũ nabaḷũ nathī} — \textquotedblleft{}This bone is not weak.\textquotedblright{}
\end{quote}

\begin{cousinweb}
\textbf{\gu{હાડકું}} (\emph{hāḍkũ}) is built on \emph{hāḍ}, \textquotedblleft{}bone\textquotedblright{}, with the suffix \emph{-kũ}. CDIAL lists \emph{hāḍ} under the lexical Sanskrit \emph{haḍḍa}, \textquotedblleft{}bone\textquotedblright{} (CDIAL 13952).
\end{cousinweb}

\begin{grammarlens}[title={What you've built}]
One more word for everyday life at A2.
\end{grammarlens}

\subsection*{Your turn}
\begin{itemize}
\raggedright
  \item \emph{Say it:} \emph{hāḍkũ}
  \item \emph{Say it:} \emph{hāḍkũ}, once more
  \item \emph{Say it:} \emph{ā hāḍkũ nabaḷũ nathī} — This bone is not weak.
  \item \emph{Recall:} say the Gujarati for injury, then the Gujarati for wound, then say \emph{hāḍkũ} again
\end{itemize}

\begin{tcolorbox}[breakable,colback=teal!4,colframe=teal!35!black,title={Before you move on}]
What does \gu{હાડકું} mean? (\textquotedblleft{}Bone\textquotedblright{}.) Say it once more.
\end{tcolorbox}
\section[\texorpdfstring{\gu{મચકોડાવું} (mackoḍāvũ)}{મચકોડાવું (mackoḍāvũ)}]{\gu{મચકોડાવું} (mackoḍāvũ) — to sprain, to get sprained}

\label{lesson:GU-C227-mackodavu}



\begin{quote}
Before the new one: say the Gujarati for wound, then the Gujarati for bone.
\end{quote}

\subsection*{You'll want to know: \gu{મચકોડાવું}}
\textbf{\gu{મચકોડાવું}} — \emph{mackoḍāvũ} — \textquotedblleft{}to sprain, to get sprained\textquotedblright{}.

\begin{quote}
\textbf{\gu{શું} \gu{તમારો} \gu{પગ} \gu{મચકોડાયો}?} — \emph{śũ tamāro pag mackoḍāyo?} — \textquotedblleft{}Did you sprain your foot?\textquotedblright{}
\end{quote}

Word for word, did your foot get sprained.

\begin{cousinweb}
The origin of \textbf{\gu{મચકોડાવું}} (\emph{mackoḍāvũ}) is not settled: CDIAL lists no Gujarati form like it. It is the \emph{-ā-} form, of a thing that happens to you, beside \emph{mackoḍvũ}, \textquotedblleft{}to twist\textquotedblright{}.
\end{cousinweb}

\begin{grammarlens}[title={What you've built}]
One more word for everyday life at A2.
\end{grammarlens}

\subsection*{Your turn}
\begin{itemize}
\raggedright
  \item \emph{Say it:} \emph{mackoḍāvũ}
  \item \emph{Say it:} \emph{mackoḍāvũ}, once more
  \item \emph{Say it:} \emph{śũ tamāro pag mackoḍāyo?} — Did you sprain your foot?
  \item \emph{Recall:} say the Gujarati for wound, then the Gujarati for bone, then say \emph{mackoḍāvũ} again
\end{itemize}

\begin{tcolorbox}[breakable,colback=teal!4,colframe=teal!35!black,title={Before you move on}]
What does \gu{મચકોડાવું} mean? (\textquotedblleft{}To sprain, to get sprained\textquotedblright{}.) Say it once more.
\end{tcolorbox}
\section[\texorpdfstring{\gu{લંગડાવું} (langḍāvũ)}{લંગડાવું (langḍāvũ)}]{\gu{લંગડાવું} (langḍāvũ) — to limp}

\label{lesson:GU-C227-langdavu}



\begin{quote}
Before the new one: say the Gujarati for bone, then the Gujarati for to sprain.
\end{quote}

\subsection*{You'll want to know: \gu{લંગડાવું}}
\textbf{\gu{લંગડાવું}} — \emph{langḍāvũ} — \textquotedblleft{}to limp\textquotedblright{}.

\begin{quote}
\textbf{\gu{તે} \gu{હવે} \gu{નથી} \gu{લંગડાતો}.} — \emph{te have nathī langḍāto} — \textquotedblleft{}He no longer limps.\textquotedblright{}
\end{quote}

\begin{cousinweb}
\textbf{\gu{લંગડાવું}} (\emph{langḍāvũ}) is a tadbhava. CDIAL lists Gujarati \emph{lãgṛũ}, \textquotedblleft{}lame\textquotedblright{}, and \emph{lãgṛāvũ} under Sanskrit \emph{laṅga}, \textquotedblleft{}lame\textquotedblright{} (CDIAL 10877).
\end{cousinweb}

\begin{grammarlens}[title={What you've built}]
One more word for everyday life at A2.
\end{grammarlens}

\subsection*{Your turn}
\begin{itemize}
\raggedright
  \item \emph{Say it:} \emph{langḍāvũ}
  \item \emph{Say it:} \emph{langḍāvũ}, once more
  \item \emph{Say it:} \emph{te have nathī langḍāto} — He no longer limps.
  \item \emph{Recall:} say the Gujarati for bone, then the Gujarati for to sprain, then say \emph{langḍāvũ} again
\end{itemize}

\begin{tcolorbox}[breakable,colback=teal!4,colframe=teal!35!black,title={Before you move on}]
What does \gu{લંગડાવું} mean? (\textquotedblleft{}To limp\textquotedblright{}.) Say it once more.
\end{tcolorbox}
